(1) 由界面切向电磁场连续条件，p偏振电场反射系数可写为 \(r_p=(n_2\cos\theta_i-n_1\cos\theta_t)/(n_2\cos\theta_i+n_1\cos\theta_t)\)，其反射率 \(R_p=|r_p|^2\)。折射角由 \(n_1\sin\theta_i=n_2\sin\theta_t\) 确定，在可传播透射区为实数。
(2) Brewster消反射要求 \(n_2\cos\theta_B=n_1\cos\theta_t\)。联立Snell定律，可得 \(\theta_B+\theta_t=\pi/2\) 和 \(\tan\theta_B=n_2/n_1\)，因此 \(\theta_B=\arctan(n_2/n_1)\)。此时反射光与折射光夹角为直角，且无p偏振反射。
(3) s偏振反射振幅 \(r_s=(n_1\cos\theta_i-n_2\cos\theta_t)/(n_1\cos\theta_i+n_2\cos\theta_t)\)。在Brewster角处利用 \(\cos\theta_B=\sin\theta_t\) 与Snell关系，得 \(r_s=(n_1^2-n_2^2)/(n_1^2+n_2^2)\)，所以 \(R_s=[(n_1^2-n_2^2)/(n_1^2+n_2^2)]^2\)，一般不为零。
(4) 若 \(n_1=n_2\)，任意入射角都无界面反射，Brewster角只是退化描述；空气至玻璃 \(n_2>n_1\) 时Brewster角大于 \(45^\circ\)。若 \(n_1>n_2\)，临界角 \(\theta_c=\arcsin(n_2/n_1)\)，满足 \(\theta_B<\theta_c\)，故Brewster消反射位于透射可传播区域。
